% vim: set foldmethod=marker foldlevel=0:

\documentclass[a4paper]{article}
\usepackage[UKenglish]{babel}

\usepackage{preamble}

\renewcommand{\thesubsection}{Q\arabic{section}~(\roman{subsection})}

\fancyhead[L]{MA266 Assignment 3}
\title{MA266 Multilinear Algebra, Assignment 3}
\colorlet{questionbodycolor}{magenta!50}

\begin{document}

\maketitle

\setlength{\parindent}{0em}
\setlength{\parskip}{1em}

% {{{ Q14
\question{14}

\begin{questionbody}
Let $p$ be prime and consider the matrix \[
M := \begin{pmatrix}
1 & 1 \\
0 & 1
\end{pmatrix} \in \op{GL}\nolimits_2({\bb F}_p).
\] Show that $M$ is not diagonalisable. You may use (without proof) the fact that a matrix is diagonalisable if and only if it is conjugate to a diagonal matrix.
\end{questionbody}

It is sufficient to show that $\forall P \in \op{GL}_2({\bb F}_p)$, $P M P^{-1}$ is not diagonal.

Let \[
P = \begin{pmatrix} a & b \\ c & d \end{pmatrix},
\] then \[
P^{-1} = \f1{ad - bc} \begin{pmatrix} d & -b \\ -c & a \end{pmatrix}.
\]

Then we get
\begin{align*}
P M P^{-1} &= \begin{pmatrix} a & b \\ c & d \end{pmatrix}
    \begin{pmatrix} 1 & 1 \\ 0 & 1 \end{pmatrix}
    \f1{ad - bc} \begin{pmatrix} d & -b \\ -c & a \end{pmatrix} \\[0.5ex]
&= \f1{ad - bc} \begin{pmatrix}
    a & a + b \\
    c & c + d
    \end{pmatrix}
    \begin{pmatrix} d & -b \\ -c & a \end{pmatrix} \\[0.5ex]
&= \f1{ad - bc} \begin{pmatrix}
    ad - ac - bc & -ab +a^2 + ab \\
    cd - c^2 - cd & -bc + ac + ad
    \end{pmatrix} \\[0.5ex]
&= \f1{ad - bc} \begin{pmatrix}
    ad - ac - bc & a^2 \\
    -c^2 & -bc + ac + ad
    \end{pmatrix}
\end{align*}

For this to be diagonal, we would need both $a^2 = 0$ and $c^2 = 0$. This is not possible since we're working in $\bb F_p$. Therefore, no choice of $P$ makes $M$ conjugate to a diagonal matrix, therefore $M$ is not diagonalisable.

\hfill $\square$

% }}}

% {{{ Q15
\newquestion{15}

\begin{questionbody}
Let $M, N \in M_n(\bb F)$ be conjugate matrices, and let $\lambda \in {\bb F}^\times$. We write $E_{M, k}(\lambda)$ for the generalised eigenspace of $M$ of index $k$ with respect to $\lambda$.
\begin{enumerate}[(i)]
\item Write $V := {\bb F}^n$. For $k \in \N$, show that $E_{M, k}(\lambda) \ne \{0_V\}$ if and only if $\lambda$ is an eigenvalue of $M$.

\item Show that $E_{M, k}(\lambda) \subset E_{M, k+1}(\lambda)$ for all $k \in \N$.

\item For $k \in \N$, a partition $n = \ell + m$ of $n$, and matrices $X \in M_\ell(\bb F)$, $Z \in M_m(\bb F)$, show that \[
\dim_{\bb F} E_{\op{diag}(X, Z), k}(\lambda) =
\dim_{\bb F} E_{X, k}(\lambda) + \dim_{\bb F} E_{Z, k}(\lambda).
\] You may use the fact that if $n = n_1 + \dotsb + n_t$, $M = \op{diag}(M_1, \dotsc, M_t)$, and \[
v = \begin{pmatrix} \lambda_1 \\ \vdots \\ \vdots \\ \lambda_n \end{pmatrix} \in {\bb F}^n,
\ v_1 = \begin{pmatrix} \lambda_1 \\ \vdots \\ \lambda_{n_1} \end{pmatrix} \in {\bb F}^{n_1},
\ v_2 = \begin{pmatrix} \lambda_{n_1 + 1} \\ \vdots \\ \lambda_{n_1 + n_2} \end{pmatrix} \in {\bb F}^{n_2},
\ \dotsc
\] then \[
M v = \begin{pmatrix} M_1 v_1 \\ \vdots \\ \vdots \\ M_t v_t \end{pmatrix}.
\]

\item Let $\lambda \in \bb F$. Show that $\dim_{\bb F} E_{M, k}(\lambda) = \dim_{\bb F} E_{N, k}(\lambda)$ for all $k \in \N$.
\end{enumerate}
\end{questionbody}

\subsection{~} % 15.i

Answer

\subsection{~} % 15.ii

Answer

\subsection{~} % 15.iii

Answer

\subsection{~} % 15.iv

Answer

% }}}

\end{document}
